\section{Chapter 10}

\begin{verse}
Conserving the One to which spirit and body are joined\\
And are no longer separated\\
To make the fresh and subtle breath circulate (in the body)\\
Begetting the embryo.\\
To polish the secret mirror excluding every complex thought\\
To which the mind does not wear itself out\\
In relations with the others and ruling the State\\
To follow non-acting\\
The instability changeableness of fortune the opening and closing of the Gate of Heaven \\
Is worth developing the receptivity of the soul = the virtue of the feminine \\
With the essential vision that embraces every aspect the four dimensions \\
To eliminate (discursive) knowing = to seem ignorant \\
In order to reach the development\\
To create without possessing\\
To act without appropriating\\
To rise without forcing\\
This is the Way
\end{verse}

The first part of the chapter reflects the ideas of operative Taoism about the elaboration of immortality. The hint about the practice with breathing, vehicle of vital energy that, subtlized, was used as means for the secret unification of the spirit and body in the Origin. Spirit and body, \emph{hun} and \emph{p'o} (first line) -- technically: yang soul and yin soul, spiritual soul and physical soul. The former is the non-terrestrial luminous principle, the other is the being of life. One could also say: the ``being'' principle and the ``life'' principle. The privileged condition of immortality is obtained through the union of the two principles, in the sense of a perfect compenetration of the
\emph{po} by the \emph{hun}, by ``life'', by ``being'', overwhelmed and occluded from entering into the current of forms. The integration of the two principles of the One is compared to the formation of an embryo -- the embryo of an existence distant from transformation. On these practices, beyond the already cited study of H Maspero, we can see the text also translated into Italian, in the \emph{Secret of the Golden Flower}.

The starting point of such practices is also sketched out: detachment from the environment, the simplification and the calm of the mind, in order to conserve one's own completely pure and integral nature. In \textbf{Chuang Tzu}, the ``abstinence of the heart'':


\begin{quotex}
Do not listen with the ear nor with the heart, but only with the spirit. To block the way of the senses, to keep the mirror of the heart pure. To keep oneself empty. From the outside, do not let things that no longer have names penetrate inside (impressions liberated from mental translations). (IV, 1)
\end{quotex}


\begin{quotex}
The transcendent Man exercises his intelligence only like a mirror: he knows without following attraction or repulsion, without leaving any traces. In such a way he is superior to everything and impartial in the face of it. (VII, 6)
\end{quotex}


\begin{quotex}
The heart of the True Man, completely calm, is like a mirror that reflects Heaven, Earth and all beings. (XIII, 1)
\end{quotex}


All this, by reconnecting to Lao Tzu's ``polishing the secret mirror''.

The successive operative developments, liking joining the ``decoagulant'' of body and spirit, are attested also in Lieh Tzu (II,3; IV, 6) in these terms:

\begin{quotex}
The spirit is condensed little by little as the body is dissolved'' (one ``etherizes'', one ``rarifies'').
\end{quotex}


We note the correspondence with the formula \emph{solve et coagula} of the Western Hermetic-alchemical tradition, which is known also in the formula ''Flesh out the spiritual, spiritualize the body, make fixed the volatile and volatilize the fixed'' (Part II of \emph{The Hermetic Tradition}, \#16)

The second part of the chapter outlines the ``non-acting'' behavior of the True Man.

\paragraph{Chinese text and literal translation}


\textbf{Chapter 10 (\begin{CJK*}{UTF8}{bsmi}第十章\end{CJK*})}

\columnratio{.4}
\begin{paracol}{2}
\begin{verse}
\begin{CJK*}{UTF8}{bsmi}
載營魄抱一，\\
能無離乎？ \\
專氣致柔，\\
能嬰兒乎？\\
滌除玄覽，\\
能無疵乎？\\
愛民治國，\\
能無知乎？\\
天門開闔，\\
能為雌乎？\\
明白四達，\\
能無爲乎？\\
生之畜之，\\
生而不有，\\
為而不恃，\\
長而不宰，\\
是謂玄德 。
\end{CJK*}
\end{verse}
\switchcolumn
\begin{verse}
If soul have been keeping on not being felt by the body,\\
How can the body and the soul not part?\\
Trying lungs to show they are delicate,\\
How can they be like babies?\\
Clean the mirror of ultimate source,\\
How can they not found their defect?\\
To love people and to govern the country,\\
How not to appoint the Jih/Zhi/wise people?\\
The gate of the heaven shuts and opens alternately and frequently,\\
How to act like females?\\
Have known all the causes and effects,\\
How to keep doing only the backlog of things?\\
Quicken them, feed them,\\
Produce and do not own.\\
Act and do not claim.\\
Raise and do not rule.\\
This is called profound virtue.
\end{verse}
\end{paracol}

\flrightit{Posted on 2020-10-04 by Lao Tzu}

